\section{The Multi-Robot Interface Problem}
\label{sec:background}

\subsection{Coupled Execution}

A multi-robot system contains a team $\mathcal R=\{r_1,\ldots,r_n\}$ whose
objectives, observations, actions, or resources are coupled
\cite{gerkey2004mrta,street2023formal,gielis2022communications}.  The coupling
may arise from a shared mission, collision
constraints, task precedence, limited resources, or communication through a
time-varying graph.  Consequently, local correctness is generally insufficient:
a decision that is safe for one robot under an old peer state can violate the
joint constraint under the current team state.

Let $x_i$, $u_i$, and $w_i$ denote the state, input, and disturbance of robot
$i$.  A general model is
\begin{equation}
  \dot x_i=f_i(x_i,u_i,w_i),
  \qquad
  g_{ij}(x_i,x_j)\ge 0,\qquad \{i,j\}\in\mathcal E_c,
  \label{eq:coupled-dynamics}
\end{equation}
where $\mathcal E_c$ indexes coupled team constraints.  For collision avoidance,
for example, $g_{ij}=\lVert p_i-p_j\rVert^2-d_{\min}^2$; other instances encode
formation, visibility, or other coupled state constraints
\cite{sharf2021dynamical,chen2019invariance,wang2017barrier}.  The local dynamics
in \cref{eq:coupled-dynamics} can be valid while the joint constraint is evaluated
with an invalid peer state.  This is the mechanical reason that horizontal
interfaces require explicit timing and uncertainty assumptions.

A typical command passes through task interpretation, allocation, prediction,
motion or action planning, distributed communication, runtime validation, and
physical execution.  These functions need not occupy separate processes or
robots, but each transformation changes the object being reasoned about.  A
language instruction becomes a task model; a task model becomes an assignment;
an assignment becomes trajectories or actions; and distributed observations
become inputs to safety-critical control.  Evidence that is meaningful at one
stage may be insufficient at the next.

\subsection{What a Contract Protects}

An assume--guarantee contract separates an environment assumption $A$ from a
component guarantee $G$ \cite{benveniste2007contracts}.  A contract theory may
also define satisfaction, compatibility, refinement, composition, and quotient
for a chosen behavioral model.  Robotics uses this idea at several levels: a
mission contract can describe temporal goals and contexts; a component contract
can specify software behavior; a dynamical contract can constrain input and
output trajectories; and a runtime-assurance specification can govern transfer
of control authority.  These contract objects are related, but they are not
interchangeable.

Here, a \emph{contract} states obligations, \emph{verification} checks a stated
property, and \emph{assurance} is the broader argument and evidence supporting
confidence in system behavior.  A contract contributes to assurance only when
its assumptions, checking mechanism, and evidence are explicit.

This distinction matters because an interface schema is not automatically a
contract.  Names, types, timestamps, and version fields become contractual only
when they participate in explicit assumptions or guarantees, are checked by a
defined validator, and have a specified consequence when the check fails.
Likewise, a runtime monitor observes only the properties encoded in its
specification; it cannot infer an omitted frame, validity horizon, or recovery
policy.  The survey therefore considers both the stated property and the
evidence used to support it.

\subsection{Evidence Boundary}

Because every guarantee depends on its assumptions, the type of supporting
evidence matters.  The survey distinguishes three forms.  \emph{Capability evidence}
shows that a system performs a task.  \emph{Assurance evidence} supports a
specified property under declared assumptions, for example through proof,
synthesis, monitoring, or testing.  \emph{Deployment evidence} measures behavior
on a stated platform, configuration, workload, or fault condition.  A planning
benchmark is therefore not evidence of network-timing robustness, and a formal
guarantee under a dynamics model is not evidence that an implementation uses
the correct frame, clock, or software version.  The survey compares approaches
within their reported evidence boundaries and does not treat all successful
demonstrations as equivalent.  This distinction is especially important for
coordinated-autonomy studies, where a mission-level probability or reliability
estimate can be valid for its model and data without certifying every runtime
interface used to execute the mission
\cite{chaki2013probabilistic,porter2013reliability}.
